\chapter{Технологическая часть}

\section{Выбор средств реализации}

В качестве системы управления базами данных была выбрана PostgreSQL~\cite{postgresql}.
Данная СУБД поддерживает реляционную модель данных, механизмы
ссылочной целостности, транзакции, триггеры и индексы различных типов.

PostgreSQL предоставляет встроенные средства оптимизации SQL-запросов,
поддерживает расширяемую систему типов данных и язык процедурного
программирования PL/pgSQL.

Для кэширования часто используемых данных был выбран Redis~\cite{redis}.
Использование Redis позволяет уменьшить количество обращений к основной
базе данных и сократить время выполнения запросов.

\section{Реализация базы данных}

База данных была реализована в PostgreSQL в соответствии с разработанной
ранее схемой.

Основными сущностями системы являются:

\begin{itemize}
	\item пользователи;
	\item игры;
	\item модификации;
	\item версии модификаций;
	\item комментарии;
	\item категории;
	\item пользовательские подборки.
\end{itemize}

Для хранения данных использовались взаимосвязанные таблицы.
Связи между таблицами реализованы с помощью внешних ключей.
В качестве идентификаторов использовались первичные ключи.
Для обеспечения ссылочной целостности были созданы внешние ключи.
Полный SQL-скрипт создания таблиц с ограничениями приведен в листинге~\ref{app:sql}.

Для хранения статуса версии модификации использовался перечислимый тип представленный в листинге~\ref{lst:enum}.

\newpage
\begin{lstlisting}[
	language=SQL,
	caption={Перечислимый тип статуса},
	label={lst:enum}
	]
	CREATE TYPE mod_status AS ENUM (
	'pending',
	'approved',
	'rejected'
	);
\end{lstlisting}

Использование перечислимого типа позволяет ограничить набор допустимых
значений поля status и предотвратить запись некорректных данных.

В результате была реализована реляционная структура базы данных,
обеспечивающая хранение взаимосвязанных данных предметной области.

\section{Ограничения целостности}

Для обеспечения корректности данных были реализованы ограничения
целостности.

В базе данных использовались следующие типы ограничений:

\begin{itemize}
	\item PRIMARY KEY;
	\item FOREIGN KEY;
	\item UNIQUE;
	\item CHECK;
	\item NOT NULL.
\end{itemize}

Ограничения предотвращают добавление некорректных данных и обеспечивают
согласованность связанных таблиц.

Для проверки допустимых значений использовались CHECK-ограничения.

Пример ограничения представлен в
листинге~\ref{lst:check_constraint}.

\begin{lstlisting}[
	language=SQL,
	caption={CHECK-ограничение},
	label={lst:check_constraint}
	]
	downloads INTEGER NOT NULL
	CHECK(downloads >= 0)
\end{lstlisting}

Использование ограничений позволяет поддерживать корректность данных
непосредственно на уровне СУБД независимо от клиентского приложения.

\section{Реализация триггеров}

В системе были реализованы триггеры для автоматического контроля
корректности данных.
Согласно требованиям предметной области, версия модификации не может
получить статус <<approved>> без прикрепленных файлов.
Для реализации данного ограничения была создана триггерная функция.

Фрагмент функции представлен в
листинге~\ref{lst:trigger_function}.

\begin{lstlisting}[
	language=SQL,
	caption={Триггерная функция},
	label={lst:trigger_function}
	]
	CREATE OR REPLACE FUNCTION
	check_mod_version_files()
	RETURNS TRIGGER AS $$
	BEGIN
		IF NEW.status = 'approved' THEN
			IF NOT EXISTS (
				SELECT 1
				FROM mod_files
				WHERE mod_version_id = NEW.id
				) THEN
				RAISE EXCEPTION
				'Version must contain at least one file';
			END IF;
		END IF;
	
		RETURN NEW;
	END;
	$$ LANGUAGE plpgsql;
\end{lstlisting}

После создания функции триггер был подключен к таблице mod\_versions.

\begin{lstlisting}[
	language=SQL,
	caption={Создание триггера},
	label={lst:create_trigger}
	]
	CREATE TRIGGER
		mod_version_validation_trigger
		BEFORE UPDATE ON mod_versions
		FOR EACH ROW
		EXECUTE FUNCTION
			check_mod_version_files();
\end{lstlisting}

Использование триггеров позволяет проверять корректность
данных на уровне базы данных и предотвращать нарушение бизнес-логики.

\section{Ролевая модель}

Для разграничения доступа к объектам базы данных была реализована
ролевая модель.

В системе используются следующие роли:

\begin{itemize}
	\item mw\_guest;
	\item mw\_user;
	\item mw\_admin.
\end{itemize}

Роль mw\_guest обладает правами только на чтение данных.

Роль mw\_user может добавлять и изменять данные, связанные с
пользовательским контентом.

Роль mw\_admin обладает полным доступом ко всем объектам базы данных.

Использование ролевой модели позволяет ограничить доступ пользователей
к данным и уменьшить вероятность несанкционированного изменения
информации.

\section{Кэширование}

Для уменьшения нагрузки на PostgreSQL использовался Redis.

В Redis сохранялись результаты наиболее часто выполняемых запросов:

\begin{itemize}
	\item популярные модификации;
	\item результаты поиска;
	\item информация о пользователях.
\end{itemize}

Использование кэширования позволяет сократить количество обращений к
основной базе данных и уменьшить время ответа системы при выполнении большого количества
однотипных запросов, результаты которых изменяются редко.

\section{Тестирование}

Для проверки корректности работы базы данных было проведено
тестирование ограничений целостности, триггеров и индексов.

В процессе тестирования проверялись:

\begin{itemize}
	\item ограничения внешних ключей;
	\item ограничения уникальности;
	\item CHECK-ограничения;
	\item работа триггеров;
	\item корректность работы индексов.
\end{itemize}

Результаты тестирования триггера представлены в
таблице~\ref{tab:trigger_test}.

\begin{table}[h]
	\centering
	\caption{Результаты тестирования триггера}
	\label{tab:trigger_test}
	\begin{tabular}{|p{6cm}|p{6cm}|}
		\hline
		Тестовый сценарий & Результат \\
		\hline
		Изменение статуса версии на approved при наличии файлов &
		Обновление выполнено успешно \\
		\hline
		Изменение статуса версии на pending &
		Обновление выполнено успешно \\
		\hline
		Изменение статуса версии на rejected &
		Обновление выполнено успешно \\
		\hline
		Изменение статуса версии на approved при отсутствии файлов &
		Сгенерирована ошибка \\
		\hline
	\end{tabular}
\end{table}

Также проводилось тестирование ограничений целостности и индексов.

Результаты тестирования показали корректную работу ограничений,
триггеров, индексов и механизмов обеспечения целостности данных.

В ходе тестирования не было обнаружено нарушений ссылочной целостности
или ошибок выполнения триггеров.

\section{Вывод}

В ходе технологической части была реализована база данных платформы
игровых модификаций с использованием PostgreSQL.Была разработана реляционная структура хранения данных, реализованы
ограничения целостности, триггеры, индексы и ролевая модель доступа. Использование ограничений и триггеров позволило обеспечить контроль
корректности данных непосредственно на уровне СУБД. Реализация ролевой модели обеспечила разграничение доступа пользователей
к объектам базы данных. Дополнительно был реализован механизм кэширования, позволяющий снизить
нагрузку на PostgreSQL при выполнении часто повторяющихся запросов. Проведенное тестирование подтвердило корректность работы разработанной
базы данных и реализованных механизмов обеспечения целостности данных.